\documentclass[11pt]{article}
\usepackage{amsmath,amssymb,mathtools}
\usepackage{fontspec,unicode-math}
\setmainfont{Latin Modern Roman}
\setmathfont{Latin Modern Math}
\usepackage{xurl,hyperref}
\newcommand{\sref}[2]{\hyperlink{source:#1:#2}{[S#2]}}
\newcommand{\cataloguescope}[1]{\paragraph{Mathematical statement.}#1}
\begin{document}
\section{There is no n so that the only integer m with phi(n) = phi(m) is m=n}\label{statement}

\subsection{Definitions and mathematical statement}

\paragraph{Shared notation.}$\mathbb N=\{1,2,\ldots\}$, and $\mathbb Z,\mathbb Q,\mathbb R,\mathbb C$ denote the integers, rationals, reals and complex numbers. Any other conventions are specified locally.


For $n\in\mathbb N$, Euler's totient is
\[\varphi(n)=\#\{a\in\mathbb N:1\le a\le n,\ \gcd(a,n)=1\},\]
so $\varphi(1)=1$.

\paragraph{Statement recorded in the supplied source.}
For every $n\in\mathbb N$, does there exist $m\in\mathbb N\setminus\{n\}$ with $\varphi(m)=\varphi(n)$? Equivalently, can any value attained by the totient function have exactly one positive-integer preimage? \sref{statement}{2}

\subsection{Short English statement}
Does every positive integer share its totient value with another positive integer?
\subsection{Sources}
\begin{description}
\item[\hypertarget{source:statement:1}{S1}] ULAM.AI and the UnsolvedMath Contributors, \emph{UnsolvedMath: A Curated Collection of Open Mathematics Problems}, catalogue entry. CC BY 4.0. \url{https://huggingface.co/datasets/ulamai/UnsolvedMath}.
\item[\hypertarget{source:statement:2}{S2}] Florentin Smarandache, \emph{On Carmichaël\textquotesingle s conjecture} (2007 arXiv posting), Introduction and Section4, explicit quantification of Carmichael’s conjecture, pages1 and3. \url{https://arxiv.org/abs/0704.2453}.
\end{description}
\label{end:statement}
\end{document}
